% --- Tabelas para SPACEMOBILE-005 ---
\begin{table}[H]
    \centering
    \caption{Elementos Orbitais do Secundário - PEGASUS DEB (SSC 24980)}
    \label{tab:sec_61046_24980}
    \small
    \begin{tabular}{lc}
        \toprule
        \textbf{Parâmetro} & \textbf{Valor} \\
        \midrule
        Semieixo Maior ($a$) & 6982,057\text{ km} \\
        Excentricidade ($e$) & 0,01547 \\
        Inclinação ($i$) & 81,58^\circ \\
        RAAN ($\Omega$) & 305,27^\circ \\
        Argumento do Perigeu ($\omega$) & 5,26^\circ \\
        Anomalia Média ($M$) & 355,02^\circ \\
        Movimento Médio ($n$) & 14,88\text{ rev/dia} \\
        Período Orbital ($T$) & 96,77\text{ min} \\
        BSTAR & $1,94450\times 10^{-3} \text{ ER}^{-1}$ \\
        \bottomrule
    \end{tabular}
\end{table}

\begin{table}[H]
    \centering
    \caption{Elementos Orbitais do Secundário - HORUS 1 (SSC 55690)}
    \label{tab:sec_61046_55690}
    \small
    \begin{tabular}{lc}
        \toprule
        \textbf{Parâmetro} & \textbf{Valor} \\
        \midrule
        Semieixo Maior ($a$) & 6874,077\text{ km} \\
        Excentricidade ($e$) & 0,00019 \\
        Inclinação ($i$) & 97,36^\circ \\
        RAAN ($\Omega$) & 128,30^\circ \\
        Argumento do Perigeu ($\omega$) & 161,35^\circ \\
        Anomalia Média ($M$) & 198,78^\circ \\
        Movimento Médio ($n$) & 15,23\text{ rev/dia} \\
        Período Orbital ($T$) & 94,53\text{ min} \\
        BSTAR & $2,40280\times 10^{-4} \text{ ER}^{-1}$ \\
        \bottomrule
    \end{tabular}
\end{table}

\begin{table}[H]
    \centering
    \caption{Elementos Orbitais do Secundário - TPA-1 (SSC 64533)}
    \label{tab:sec_61046_64533}
    \small
    \begin{tabular}{lc}
        \toprule
        \textbf{Parâmetro} & \textbf{Valor} \\
        \midrule
        Semieixo Maior ($a$) & 6882,710\text{ km} \\
        Excentricidade ($e$) & 0,00042 \\
        Inclinação ($i$) & 97,46^\circ \\
        RAAN ($\Omega$) & 167,34^\circ \\
        Argumento do Perigeu ($\omega$) & 226,33^\circ \\
        Anomalia Média ($M$) & 133,76^\circ \\
        Movimento Médio ($n$) & 15,20\text{ rev/dia} \\
        Período Orbital ($T$) & 94,71\text{ min} \\
        BSTAR & $1,95200\times 10^{-4} \text{ ER}^{-1}$ \\
        \bottomrule
    \end{tabular}
\end{table}

% --- Tabelas para KUIPER-00084 ---
\begin{table}[H]
    \centering
    \caption{Elementos Orbitais do Secundário - COSMOS 2251 DEB (SSC 33825)}
    \label{tab:sec_64830_33825}
    \small
    \begin{tabular}{lc}
        \toprule
        \textbf{Parâmetro} & \textbf{Valor} \\
        \midrule
        Semieixo Maior ($a$) & 7035,740\text{ km} \\
        Excentricidade ($e$) & 0,00816 \\
        Inclinação ($i$) & 73,98^\circ \\
        RAAN ($\Omega$) & 252,49^\circ \\
        Argumento do Perigeu ($\omega$) & 203,81^\circ \\
        Anomalia Média ($M$) & 279,71^\circ \\
        Movimento Médio ($n$) & 14,71\text{ rev/dia} \\
        Período Orbital ($T$) & 97,89\text{ min} \\
        BSTAR & $9,15790\times 10^{-4} \text{ ER}^{-1}$ \\
        \bottomrule
    \end{tabular}
\end{table}

\begin{table}[H]
    \centering
    \caption{Elementos Orbitais do Secundário - COSMOS 2251 DEB (SSC 36045)}
    \label{tab:sec_64830_36045}
    \small
    \begin{tabular}{lc}
        \toprule
        \textbf{Parâmetro} & \textbf{Valor} \\
        \midrule
        Semieixo Maior ($a$) & 7032,829\text{ km} \\
        Excentricidade ($e$) & 0,00475 \\
        Inclinação ($i$) & 73,76^\circ \\
        RAAN ($\Omega$) & 282,33^\circ \\
        Argumento do Perigeu ($\omega$) & 284,39^\circ \\
        Anomalia Média ($M$) & 75,20^\circ \\
        Movimento Médio ($n$) & 14,72\text{ rev/dia} \\
        Período Orbital ($T$) & 97,83\text{ min} \\
        BSTAR & $1,21120\times 10^{-3} \text{ ER}^{-1}$ \\
        \bottomrule
    \end{tabular}
\end{table}

\begin{table}[H]
    \centering
    \caption{Elementos Orbitais do Secundário - KUIPER-00041 (SSC 64502)}
    \label{tab:sec_64830_64502}
    \small
    \begin{tabular}{lc}
        \toprule
        \textbf{Parâmetro} & \textbf{Valor} \\
        \midrule
        Semieixo Maior ($a$) & 7006,487\text{ km} \\
        Excentricidade ($e$) & 0,00102 \\
        Inclinação ($i$) & 51,90^\circ \\
        RAAN ($\Omega$) & 323,23^\circ \\
        Argumento do Perigeu ($\omega$) & 203,49^\circ \\
        Anomalia Média ($M$) & 156,56^\circ \\
        Movimento Médio ($n$) & 14,80\text{ rev/dia} \\
        Período Orbital ($T$) & 97,28\text{ min} \\
        BSTAR & $5,96080\times 10^{-4} \text{ ER}^{-1}$ \\
        \bottomrule
    \end{tabular}
\end{table}

\begin{table}[H]
    \centering
    \caption{Elementos Orbitais do Secundário - KUIPER-00152 (SSC 65772)}
    \label{tab:sec_64830_65772}
    \small
    \begin{tabular}{lc}
        \toprule
        \textbf{Parâmetro} & \textbf{Valor} \\
        \midrule
        Semieixo Maior ($a$) & 7008,706\text{ km} \\
        Excentricidade ($e$) & 0,00007 \\
        Inclinação ($i$) & 51,90^\circ \\
        RAAN ($\Omega$) & 107,73^\circ \\
        Argumento do Perigeu ($\omega$) & 207,19^\circ \\
        Anomalia Média ($M$) & 18,78^\circ \\
        Movimento Médio ($n$) & 14,80\text{ rev/dia} \\
        Período Orbital ($T$) & 97,32\text{ min} \\
        BSTAR & $-2,67220\times 10^{-3} \text{ ER}^{-1}$ \\
        \bottomrule
    \end{tabular}
\end{table}

% --- Tabelas para KUIPER-00155 ---
\begin{table}[H]
    \centering
    \caption{Elementos Orbitais do Secundário - AEROCUBE 8D (SSC 41852)}
    \label{tab:sec_65774_41852}
    \small
    \begin{tabular}{lc}
        \toprule
        \textbf{Parâmetro} & \textbf{Valor} \\
        \midrule
        Semieixo Maior ($a$) & 6876,393\text{ km} \\
        Excentricidade ($e$) & 0,00035 \\
        Inclinação ($i$) & 98,08^\circ \\
        RAAN ($\Omega$) & 272,42^\circ \\
        Argumento do Perigeu ($\omega$) & 159,54^\circ \\
        Anomalia Média ($M$) & 200,60^\circ \\
        Movimento Médio ($n$) & 15,23\text{ rev/dia} \\
        Período Orbital ($T$) & 94,58\text{ min} \\
        BSTAR & $3,17670\times 10^{-4} \text{ ER}^{-1}$ \\
        \bottomrule
    \end{tabular}
\end{table}

\begin{table}[H]
    \centering
    \caption{Elementos Orbitais do Secundário - GAOFEN 11-05 (SSC 60249)}
    \label{tab:sec_65774_60249}
    \small
    \begin{tabular}{lc}
        \toprule
        \textbf{Parâmetro} & \textbf{Valor} \\
        \midrule
        Semieixo Maior ($a$) & 6873,454\text{ km} \\
        Excentricidade ($e$) & 0,00070 \\
        Inclinação ($i$) & 97,41^\circ \\
        RAAN ($\Omega$) & 131,14^\circ \\
        Argumento do Perigeu ($\omega$) & 152,17^\circ \\
        Anomalia Média ($M$) & 207,99^\circ \\
        Movimento Médio ($n$) & 15,23\text{ rev/dia} \\
        Período Orbital ($T$) & 94,52\text{ min} \\
        BSTAR & $2,33070\times 10^{-4} \text{ ER}^{-1}$ \\
        \bottomrule
    \end{tabular}
\end{table}

